\documentclass[../../main.tex]{subfiles}
\begin{document}
\section{Open targets and roadmap}
\label{sec:open}

The consolidation leaves a sharper, smaller set of open targets than the original program. We
state them in decreasing order of strength.

\begin{question}[Operator-to-trace upgrade; all-cut Carleson]
\label{q:upgrade}
Prove Assumption~\ref{ass:all-cut-carleson}: upgrade the unconditional per-direction Carleson
estimate (Corollary~\ref{cor:per-direction}) to the trace scale, uniformly over balanced cuts,
with damping coefficient $\alpha<1$. By Theorem~\ref{thm:intro-all-cut} this alone implies KLS.
The projection-test ceiling of Section~\ref{sec:qcts} shows that radial and projection
information caps at $\log n$; the upgrade must use the cut-specific structure. The first
decision point is Program~\ref{prog:product-test}.
\end{question}

\begin{question}[Weighted excess propagation with rate]
\label{q:weighted}
Prove \eqref{eq:intro-weighted-excess}: for balanced near-Cheeger cuts,
\[
  \E\int_0^{T\wedge\tau_\eta}e_t(E)\,\bigl(1+\norm{A_t}_\op\bigr)^{5/2}\dd t
  \ \le\ C\bigl(Te_0(E)+T^{1+\gamma}\bigr),\qquad T\le T_0 .
\]
By Remark~\ref{rem:weight-explains-rate} this is the weakest excess statement statically
consistent with the two-tail mode, with exponent calibrated on
Proposition~\ref{prop:two-tail}; the bootstrap of Theorem~\ref{thm:bootstrap} resolves the
unweighted analogue but the weighted version requires, in addition, second-moment control of
the perimeter martingale against the covariance weight, which is open.
\end{question}

\begin{question}[Weighted stable Stein trace]
\label{q:stein-weighted}
Prove \eqref{eq:intro-weighted-stein} for balanced near-Cheeger cuts, with damping coefficient
$\beta<\tfrac12$, by the mechanism of Section~\ref{sec:jacobi}: quadratic-chaos control of the
interior Dirichlet quantity, stability control of mean-zero boundary modes via \eqref{eq:stability},
and the constant mode via the Schur energy \eqref{eq:Jsharp-def} and splitting. The chaining must be
uniform along the localization and at the covariance-weighted scale; by
Proposition~\ref{prop:two-tail} no slice-wise shortcut exists.
\end{question}

\begin{remark}[The three open trace-upgrade targets are one problem]
\label{rem:trace-upgrade-unification}
Questions~\ref{q:upgrade} and \ref{q:stein-weighted} (its high-rank, mean-zero part) and the
adapted alignment problem~\ref{q:alignment} are the \emph{same} operator-to-trace difficulty
in three coordinate systems. Analytically, the gap is exactly the trace of the occupation
operator $M:=\E\int_0^\infty s_t G_t^2\dd t$: Corollary~\ref{cor:per-direction} is the
statement $M\preceq R_0\preceq I$ (the diagonal/quadratic-form level), and the trace target is
$\Tr(M)$, whose naive bound $\sum_i(R_0)_{ii}=\Tr R_0\le n$ is the obstruction. Geometrically,
the high-rank mean-zero residue of Question~\ref{q:stein-weighted} is the same sum, and its
constant mode is Question~\ref{q:splitting}. Combinatorially, summing the per-coordinate
budgets of Theorem~\ref{thm:budget} over a product cut is again $\Tr R_0\le n$, i.e.\
Question~\ref{q:alignment}. The correct target is therefore \emph{not} a static
effective-rank bound --- on products $\Tr(M)$ genuinely can be of order $n$, yet the all-cut
estimate still holds on the polylog window (Theorem~\ref{thm:V2-window}) because the $n$
coordinate budgets are spent at \emph{temporally disjoint} times --- but an
\emph{occupation-density} bound: the source-occupation measure on time$\,\times\,$direction
should be absolutely continuous with bounded density on the early balanced window, so that
only $O(1)$ directions are simultaneously active. This is the form in which
Question~\ref{q:alignment} should be integrated against time.
\end{remark}

\begin{question}[Extremality tames the covariance process]
\label{q:taming}
For every $\kappa>0$, do there exist $\eps,T_0>0$ such that every isotropic log-concave $\mu$
on $\R^n$ with $h_\mu\le(1+\eps)\hstar_n$ satisfies
$h_\mu\,\Xi_{T_0}(\mu)\le\kappa\,T_0$? By Theorem~\ref{thm:bootstrap} an affirmative answer
gives relative-scale excess propagation for near-worst measures, while by
Proposition~\ref{prop:ceiling} the corresponding statement without the near-worstness
hypothesis is KLS-equivalent and hence not a legitimate target. Heuristically, covariance
inflation under localization is driven by third-moment anisotropy, while the splitting
philosophy of Section~\ref{sec:jacobi} predicts that near-worst measures are approximately
split along their dangerous directions, where the variance processes are one-dimensional and
tame (Proposition~\ref{prop:products}(ii)); making this rigorous would couple the splitting
analysis to the covariance SDE.
\end{question}

\begin{remark}[The naive splitting--covariance coupling faces a type mismatch]
\label{rem:taming-type-mismatch}
The coupling suggested above runs into a
structural obstruction: $\Xi_T$ is a \emph{cut-free} functional of the covariance process,
whereas the splitting structure is \emph{cut-indexed}, and the split direction of
Proposition~\ref{prop:exact-splitting} is precisely the inflating (log-affine, exponential)
direction along which $\lmax(A_t)$ grows --- so the splitting philosophy controls the
\emph{source} of the cut-aware Stein estimate (Questions~\ref{q:stein-weighted},
\ref{q:splitting}), not the cut-free $\Xi_T$. The residual non-circular content of
Question~\ref{q:taming} is therefore a \emph{near-worst-specific} improvement of the
Klartag--Lehec window precisely in the regime $\hstar_n\gg1/\log\log n$ that
Corollary~\ref{cor:loglog} leaves open.
\end{remark}

\begin{question}[Quantitative splitting / boundary Obata]
\label{q:splitting}
Prove a quantitative version of the rigidity principle of Remark~\ref{rem:obata}: for a
near-worst measure and a near-minimal balanced cut, small constant-mode curvature
$\mathfrak K_\Sigma$ forces quantitative proximity, along the normal direction, to the split
structure of Proposition~\ref{prop:exact-splitting}, with constants stable under the
localization tilt (Proposition~\ref{prop:persistent-splitting}).
\end{question}

\begin{corollary}[Residual dichotomy of the bootstrap route]
\label{cor:dichotomy}
Assume Hypothesis~\ref{hyp:KI}. Suppose the geometric route is completed with an excess demand
at absolute scale with sufficiently small constant $\kappa$ --- that is, suppose the weighted
Stein-trace mechanism can run on the supply
$\int_0^T\E[\bar e_t\one_{t<\tau_\eta}]\dd t\le Te_0+\kappa T$. Then, by
Corollary~\textup{\ref{cor:loglog}}, this supply is available for near-worst measures in every
dimension with $\hstar_n\le c\kappa T_0/(1+\log\log n)$, and the resulting contradiction
confines any failure of KLS to dimensions with $\hstar_n\ge c'/\log\log n$. Such a completion
would therefore already yield
\[
  \hstar_n\ \ge\ \frac{c'(\kappa)}{\log\log n}
  \qquad\text{for all large }n,
\]
an exponential improvement of the best known bound $\hstar_n\ge c(\log n)^{-1/2}$
\cite{Klartag2023Logarithmic}, with full KLS following from an affirmative answer to
Question~\textup{\ref{q:taming}} in the remaining regime.
\end{corollary}

\begin{proof}
Choose the numerical constant $c$ in the hypothesis small enough, depending only on the constants
in Corollary~\ref{cor:loglog} and on the completed route. If
$\hstar_n\le c\kappa T_0/(1+\log\log n)$, pick an isotropic log-concave $\mu$ in dimension $n$ with
$h_\mu\le2\hstar_n$ and a balanced near-minimizer $E$ satisfying
$\mu^+(E)\le I_\mu(\tfrac12)+o(1)=\tfrac12h_\mu+o(1)$, where the $o(1)$ is along the
near-minimizing sequence. Corollary~\ref{cor:loglog} gives the supply
$\int_0^{T_0}\E[\bar e_t\one_{t<\tau_\eta}]\dd t\le T_0e_0+\kappa T_0$ after reducing $c$.
Running the completed route on this near-minimizing sequence yields a universal lower bound
$\mu^+(E)\ge c''>0$. Letting the near-minimizer error tend to zero gives $h_\mu\ge2c''$, which
contradicts $h_\mu\le2\hstar_n$ once $n$ is in the asserted small-$\hstar_n$ regime. Hence any
failure of KLS must lie in the complementary regime $\hstar_n\gtrsim1/\log\log n$.
\end{proof}

\subsection{Roadmap}

\begin{center}
\begin{tabularx}{\textwidth}{@{}>{\raggedright\arraybackslash}p{0.27\textwidth}>{\raggedright\arraybackslash}p{0.34\textwidth}Y@{}}
\toprule
Step & Input & Output \\
\midrule
Program~\ref{prog:product-test} & One-dimensional log-concave analysis; explicit variance
dynamics & Decides whether Route A survives; either outcome redirects effort decisively. \\
\addlinespace
Question~\ref{q:weighted} & Perimeter martingale moments against the covariance weight;
bootstrap technology of Section~\ref{sec:bootstrap} & Propagation half of the weighted
package. \\
\addlinespace
Question~\ref{q:stein-weighted} & Reilly identity, stability inequality, quadratic chaos;
constant mode via Question~\ref{q:splitting} & Stability half of the weighted package; with
Question~\ref{q:weighted}, KLS by Theorem~\ref{thm:intro-weighted}. \\
\addlinespace
Question~\ref{q:taming} & Coupling of splitting structure to the covariance SDE & Relative
propagation for near-worst measures; with Corollary~\ref{cor:dichotomy}, intermediate
unconditional bounds at the $\log\log n$ scale. \\
\bottomrule
\end{tabularx}
\end{center}

\subsection{Conclusion}

This report proves no bound on $\hstar_n$. Its unconditional content is: the exact stochastic
layer (Sections~\ref{sec:mass-martingale}--\ref{sec:carleson}); the two-color Stein dictionary
and the operator-to-trace gap (Section~\ref{sec:stein-dictionary}); the static obstruction ---
quadratic chaos and the quantified two-tail configuration with its weight calibration
(Section~\ref{sec:qcts}); the covariance technology and its discharge
(Section~\ref{sec:covariance-tech}); the consumption audit, the perimeter-martingale and
excess identities, and the circularity analysis (Section~\ref{sec:excess}); the bootstrap
comparison theorem with its interface evaluation and ceiling (Section~\ref{sec:bootstrap}); the
exact model statements of the splitting mechanism (Section~\ref{sec:jacobi}); and the model
geometries (Section~\ref{sec:models}). The conditional content is concentrated in the weighted
package (Assumption~\ref{ass:weighted-package}) and the all-cut estimate
(Assumption~\ref{ass:all-cut-carleson}), with the open targets posed in
Questions~\ref{q:upgrade}--\ref{q:splitting}. The architecture rests on three structural
points: the previously postulated excess-propagation assumption is retired (unconditionally
true as consumed, hence inert); the Stein-trace estimate is stated in the covariance-weighted
form forced by the static obstruction; and the residual difficulty of propagation is
compressed into the single scalar interface $h_\mu\,\Xi_T$, evaluated to within a
$\log\log n$ factor of what the route requires.
\end{document}
